% SC2005 Ch04 — Synchronization (no \documentclass)
\chapter{Synchronization}
\label{ch:sync}

\begin{quote}
\emph{Concurrency without synchronization is chaos.} This chapter builds the
primitives that make correct concurrent programs possible.
\end{quote}

\noindent\textbf{Learning Outcomes.} After this chapter you will be able to:
\begin{enumerate}[label=(L\arabic*)]
  \item identify race conditions and define the critical-section problem;
  \item explain Peterson's algorithm and hardware atomics (TSL, CAS);
  \item use semaphores and mutexes correctly;
  \item use monitors / condition variables and solve producer--consumer;
  \item avoid common pitfalls (deadlock, missed wakeup, busy waiting).
\end{enumerate}

% ============================================================
\section{Races and the Critical Section}
\label{sec:race-cs}

A \emph{race condition} occurs when unsynchronized accesses to shared state
produce nondeterministic results. Example: \texttt{count++} is
\texttt{load--inc--store}; interleaving two increments can lose an update.

\subsection{Critical-Section Requirements}
A correct mutual-exclusion protocol guarantees:
\begin{enumerate}[label=(\roman*)]
  \item \textbf{Mutual exclusion:} at most one process in its critical section (CS).
  \item \textbf{Progress:} a process not in CS cannot block others indefinitely.
  \item \textbf{Bounded waiting:} no indefinite starvation.
\end{enumerate}

\subsection{Peterson's Algorithm (2 processes)}
For processes $P_0, P_1$ with shared \texttt{flag[2]} and \texttt{turn}:

\begin{lstlisting}[language=C,caption={Peterson's algorithm for $P_i$}]
// shared: flag[2]={0,0}; turn=0;
flag[i] = 1; turn = 1 - i;           // intend + yield
while (flag[1-i] && turn == 1-i) ;   // busy-wait
// --- critical section ---
flag[i] = 0;
// --- remainder ---
\end{lstlisting}

\begin{figure}[h!]
\centering
\begin{tikzpicture}[scale=0.74, every node/.style={font=\scriptsize},
  s/.style={draw, rounded corners, align=center, minimum width=2.6cm, minimum height=0.75cm}]
  \node[s, fill=blue!12] (a) at (0,1.6) {\texttt{flag[i]=1; turn=j}};
  \node[s, fill=orange!15] (b) at (4.2,1.6) {\texttt{while(flag[j]\&\&turn==j)};};
  \node[s, fill=green!14] (c) at (4.2,0.3) {Critical Section};
  \node[s, fill=violet!12] (d) at (0,0.3) {\texttt{flag[i]=0}};
  \draw[-{Stealth}, thick, blue!60!black] (a)--(b);
  \draw[-{Stealth}, thick, green!50!black] (b)--(c) node[midway,right,font=\tiny]{guard false};
  \draw[-{Stealth}, thick, red!60!black] (b) -- ++(0,1.0) -- ++(-1.5,0) |- (b) node[pos=0.25,left,font=\tiny]{spin};
  \draw[-{Stealth}, thick] (c)--(d); \draw[-{Stealth}, thick, dashed] (d) |- (a);
\end{tikzpicture}
\caption{Peterson's protocol flow for process $P_i$ (busy-waiting).}
\label{fig:peterson}
\end{figure}

Peterson satisfies all three requirements but uses busy waiting, is limited to
two processes, and assumes sequentially consistent memory (fails under modern
reordering without fences).

\subsection{Hardware Support}
Atomic instructions let software build synchronization without pure busy-wait
loops: \texttt{test-and-set} (TSL), \texttt{compare-and-swap} (CAS), and
\texttt{fetch-and-add} --- each completes atomically in hardware.

\begin{lstlisting}[language=C,caption={Spinlock via compare-and-swap}]
#include <stdatomic.h>
atomic_flag lock = ATOMIC_FLAG_INIT;
void acquire() { while (atomic_flag_test_and_set(&lock)) ; /* spin */ }
void release() { atomic_flag_clear(&lock); }
\end{lstlisting}

% ============================================================
\section{Semaphores and Mutexes}
\label{sec:semaphores}

A \emph{semaphore} $S$ is an integer with atomic \texttt{wait(S)} / \texttt{P}
and \texttt{signal(S)} / \texttt{V}:

\begin{center}
\texttt{wait(S): while(S$\le$0); S--\quad\quad signal(S): S++}
\end{center}
Binary semaphore ($S\in\{0,1\}$) acts as a mutex; counting semaphore tracks
available resources. Practical semaphores block instead of spinning (queue +
\texttt{sleep}/\texttt{wakeup}).

\begin{figure}[h!]
\centering
\begin{tikzpicture}[scale=0.74, every node/.style={font=\scriptsize},
  b/.style={draw, rounded corners, align=center, minimum width=2.8cm, minimum height=0.75cm}]
  \node[b, fill=blue!12] (w) at (0,1.2) {\texttt{wait(S)}: if $S>0$ decr else block};
  \node[b, fill=green!14] (cs) at (0,0.0) {Critical Section};
  \node[b, fill=orange!15] (s) at (0,-1.2) {\texttt{signal(S)}: incr; wake one waiter};
  \draw[-{Stealth}, thick, blue!60!black] (w)--(cs);
  \draw[-{Stealth}, thick, blue!60!black] (cs)--(s);
  \draw[-{Stealth}, thick, dashed, red!60!black] (s.east) -- ++(0.9,0) |- (w.east) node[pos=0.25,right,font=\tiny]{next waiter};
  \node[draw, dashed, rounded corners, fit=(w)(s), inner sep=4pt, label={[font=\tiny]above right:semaphore invariant $S\ge 0$}] {};
\end{tikzpicture}
\caption{Semaphore-based mutual exclusion with blocking (not spinning).}
\label{fig:semaphore}
\end{figure}

\subsection{Producer--Consumer with Semaphores}
Buffer size $N$: \texttt{empty} $=N$, \texttt{full} $=0$, \texttt{mutex} $=1$.
\begin{lstlisting}[language=C,caption={Bounded buffer --- semaphore solution}]
sem_t empty, full, mutex; // init N, 0, 1
void *producer(void *a) {
    for (;;) { Item x = produce();
        sem_wait(&empty); sem_wait(&mutex);
        buffer[in] = x; in = (in+1)%N;
        sem_post(&mutex); sem_post(&full);
    }
}
void *consumer(void *a) {
    for (;;) {
        sem_wait(&full); sem_wait(&mutex);
        Item x = buffer[out]; out=(out+1)%N;
        sem_post(&mutex); sem_post(&empty);
        consume(x);
    }
}
\end{lstlisting}
Ordering matters: waiting on \texttt{mutex} before \texttt{empty/full} can
deadlock (see Chapter~5).

% ============================================================
\section{Monitors and Condition Variables}
\label{sec:monitors}

A \emph{monitor} encapsulates shared data + procedures with implicit mutual
exclusion (only one thread inside at a time). \emph{Condition variables}
(\texttt{wait}, \texttt{signal}/\texttt{broadcast}) let threads sleep until a
predicate holds --- always re-check in a \texttt{while} loop (spurious wakeups).

\begin{lstlisting}[language=C,caption={Monitor-style bounded buffer (pthreads)}]
pthread_mutex_t m = PTHREAD_MUTEX_INITIALIZER;
pthread_cond_t not_full = PTHREAD_COND_INITIALIZER,
               not_empty = PTHREAD_COND_INITIALIZER;
int cnt=0, in=0, out=0; Item buf[N];
void put(Item x){ pthread_mutex_lock(&m);
    while(cnt==N) pthread_cond_wait(&not_full,&m);
    buf[in]=x; in=(in+1)%N; cnt++;
    pthread_cond_signal(&not_empty); pthread_mutex_unlock(&m);
}
Item get(){ pthread_mutex_lock(&m);
    while(cnt==0) pthread_cond_wait(&not_empty,&m);
    Item x=buf[out]; out=(out+1)%N; cnt--;
    pthread_cond_signal(&not_full); pthread_mutex_unlock(&m); return x;
}
\end{lstlisting}

\begin{figure}[h!]
\centering
\begin{tikzpicture}[scale=0.74, every node/.style={font=\scriptsize},
  b/.style={draw, rounded corners, align=center, minimum width=2.6cm, minimum height=0.75cm}]
  \node[b, fill=blue!12] (enter) at (0,1.8) {Enter monitor\\(acquire lock)};
  \node[b, fill=orange!15] (check) at (0,0.7) {Predicate?};
  \node[b, fill=green!14] (work) at (3.0,0.7) {Access shared data};
  \node[b, fill=violet!12] (wait) at (0,-0.55) {\texttt{cond\_wait}\\ (release + sleep)};
  \node[b, fill=gray!14] (exit) at (3.0,-0.55) {Signal + exit\\(release lock)};
  \draw[-{Stealth}, thick] (enter)--(check);
  \draw[-{Stealth}, thick, green!50!black] (check)--(work) node[midway,above,font=\tiny]{true};
  \draw[-{Stealth}, thick, red!60!black] (check)--(wait) node[midway,left,font=\tiny]{false};
  \draw[-{Stealth}, thick, dashed, red!60!black] (wait) |- (check);
  \draw[-{Stealth}, thick] (work)--(exit);
\end{tikzpicture}
\caption{Monitor + condition variable pattern: always re-check predicate after wakeup.}
\label{fig:monitor}
\end{figure}

% ============================================================
\begin{remark}[Spurious wakeups]
On some systems \texttt{pthread\_cond\_wait} may return without a signal.
Hence the predicate must be rechecked in a loop --- an \texttt{if} would
proceed with a false precondition.
\end{remark}

\section{Exercises}
\label{sec:ch04-ex}
\begin{exercise}
Show the lost-update race for \texttt{count++} with a concrete interleaving
of \texttt{load/inc/store} instructions.
\end{exercise}
\begin{exercise}
Prove Peterson's algorithm guarantees mutual exclusion. Which requirement fails
if \texttt{turn} is removed?
\end{exercise}
\begin{exercise}
Swap the two \texttt{sem\_wait} lines in producer (\texttt{mutex} before
\texttt{empty}). Construct a deadlock scenario.
\end{exercise}
\begin{exercise}
Implement a reader--writer lock with semaphores that gives writers priority.
Explain starvation of readers.
\end{exercise}
\begin{exercise}
Why must \texttt{pthread\_cond\_wait} be called inside a \texttt{while} loop
testing the predicate, not an \texttt{if}?
\end{exercise}
